\SourceLeaf{281}%
\(x\). Diese Coordinaten sind als in Theilen des Halbmessers ausgedrückt verstanden, und müssen demnach noch multiplicirt werden mit \(A\), wenn man sie nach ihrer Lineargrösse, oder mit \(206265''\), wenn man sie in Bogentheilen ausgedrückt verlangt.

Ein Element von \(M\) wird durch
\[
\BareRoot{(\cos y^2\mathrm{d}x^2+\mathrm{d}y^2)}
\]
oder durch \(\frac{\cos y}{\cos\psi}\mathord{.}\mathrm{d}x\) ausgedrückt, wenn man
\[
\frac{\mathrm{d}y}{\cos y\mathrm{d}x}=\operatorname{tang}\psi
\]
setzt, wo mithin \(\psi\) die Neigung des Elements gegen die Parallele mit \(N\) bedeutet. Um die Vorstellung zu fixiren, mag man sich die \(x\) von der Rechten nach der Linken, die \(y\) von unten nach oben wachsend denken, wodurch also der Sinn positiver \(\psi\) von selbst bestimmt ist.

Das wie oben mit \(m\) bezeichnete Vergrösserungsverhältniss beim Uebertragen der ellipsoidischen Fläche auf die Kugelfläche kann hier wie eine Function von \(x\) und \(y\) betrachtet werden: die Grösse des Elements von \(L\), dem jenes Element von \(M\) entspricht, wird
\[
=\frac{A\cos y}{m\cos\psi}\mathord{.}\mathrm{d}x
\]
sein, und wenn zur Abkürzung
\[
\begin{aligned}
\log\operatorname{tang}(45^0+\tfrac{1}{2}y)&=u\\
\frac{\cos y}{m}&=n
\end{aligned}
\]
gesetzt wird, wo mithin \(n\) gleichfalls Function von \(x\) und \(y\), oder was auf Eines hinausläuft, von \(x\) und \(u\) sein wird, so hat man
\[
\operatorname{tang}\psi=\frac{\mathrm{d}u}{\mathrm{d}x}
\]
und das Element von \(L\)
\[
=\frac{An}{\cos\psi}\mathord{.}\mathrm{d}x
\]
\indent Die Natur der Linie \(M\) wird also durch die Bedingung bestimmt. dass zwischen irgendwelchen bestimmten Grenzen das Integral \(\int\frac{n}{\cos\psi}\mathrm{d}x\) oder
\[
\int n\BareRoot{(1+\frac{\mathrm{d}u^2}{\mathrm{d}x^2})}\mathrm{d}x
\]
